% !TeX root = ../main.tex
\chapter{ĐÁNH GIÁ KẾT QUẢ VÀ THẢO LUẬN}
\label{ch:danhgia}

\section{Phương pháp đánh giá}
\label{sec:ppdg}

Hệ thống được đánh giá theo ba khía cạnh. Thứ nhất, về \textbf{chức năng}, các
ca kiểm thử được thực hiện cho luồng đăng nhập SSO, cơ chế cấu hình realm tự
động và xác thực đa yếu tố. Thứ hai, về \textbf{bảo mật}, các mối đe dọa được
phân tích theo mô hình STRIDE và đối chiếu với biện pháp đã triển khai. Thứ ba,
về \textbf{vận hành}, mức độ tự động hóa và khả năng tái lập của hệ thống được
đánh giá. Việc kiểm thử được tiến hành trên môi trường triển khai cục bộ (Docker
Compose) và môi trường cụm (k3s).

\section{Kết quả kiểm thử chức năng}
\label{sec:ketqua-cn}

Bảng~\ref{tab:testcases} trình bày các ca kiểm thử chức năng đã thực hiện cùng
kết quả thu được.

\begin{table}[htbp]
\centering
\caption{Kết quả kiểm thử chức năng}
\label{tab:testcases}
\small
\begin{tabular}{c p{4.4cm} p{2.4cm} p{4.0cm}}
\toprule
\textbf{Mã} & \textbf{Ca kiểm thử} & \textbf{Kết quả} & \textbf{Ghi nhận}\\
\midrule
TC1 & Đăng nhập đúng thông tin & Đạt & Nhận token và truy cập ứng dụng\\
TC2 & Đăng nhập sai mật khẩu & Đạt & Bị từ chối và có ghi log\\
TC3 & Xác thực bằng mã TOTP & Đạt & Yêu cầu nhập đúng mã\\
TC4 & Access token hết hạn & Đạt & Bị từ chối hoặc được làm mới\\
TC5 & Áp dụng cấu hình realm từ JSON & Đạt & Realm được cập nhật theo mã nguồn\\
TC6 & Chạy nhiều bản sao & Đạt & Các Pod cùng phục vụ, phiên ổn định\\
\bottomrule
\end{tabular}
\end{table}

Kết quả cho thấy các chức năng cốt lõi của hệ thống hoạt động đúng như thiết kế.
Bảng~\ref{tab:ketqua} đối chiếu kết quả này với các yêu cầu chức năng đã đặt ra
ở Chương~\ref{ch:thietke}.

\begin{table}[htbp]
\centering
\caption{Đối chiếu yêu cầu chức năng và kết quả}
\label{tab:ketqua}
\small
\begin{tabular}{c p{8.2cm} c}
\toprule
\textbf{Mã} & \textbf{Yêu cầu} & \textbf{Kết quả}\\
\midrule
FN1 & Dịch vụ định danh tập trung với Keycloak & Đạt\\
FN2 & Đăng nhập SSO vào ứng dụng demo qua OIDC & Đạt\\
FN3 & Quản lý realm/client/role/user như mã nguồn & Đạt\\
FN4 & Xác thực đa yếu tố (TOTP) & Đạt\\
FN5 & Áp dụng cấu hình tự động qua CI/CD & Đạt\\
FN6 & Ứng dụng demo hiển thị thông tin từ token & Đạt\\
\bottomrule
\end{tabular}
\end{table}

% TODO: chèn ảnh chụp màn hình minh chứng cho các ca kiểm thử.

\section{Đánh giá bảo mật}
\label{sec:stride}

Các mối đe dọa được phân tích theo mô hình STRIDE và đối chiếu với biện pháp đã
áp dụng trong hệ thống (Bảng~\ref{tab:stride}).

\begin{table}[htbp]
\centering
\caption{Phân tích mối đe dọa theo STRIDE}
\label{tab:stride}
\small
\begin{tabular}{p{2.4cm} p{5.0cm} p{5.4cm}}
\toprule
\textbf{Loại} & \textbf{Mối đe dọa} & \textbf{Biện pháp đã áp dụng}\\
\midrule
Spoofing & Giả mạo người dùng, đánh cắp phiên & MFA, thời hạn token ngắn, PKCE\\
Tampering & Sửa đổi token hoặc cấu hình & Chữ ký JWT, cấu hình quản lý qua Git\\
Repudiation & Phủ nhận hành động & Ghi log kiểm toán của Keycloak\\
Information disclosure & Lộ dữ liệu nhạy cảm & TLS, bí mật tách khỏi mã nguồn\\
Denial of service & Làm quá tải IdP & Giới hạn tần suất, chạy nhiều bản sao\\
Elevation of privilege & Leo thang đặc quyền & RBAC, đặc quyền tối thiểu, NetworkPolicy\\
\bottomrule
\end{tabular}
\end{table}

Kết quả phân tích cho thấy việc kết hợp SSO tập trung, MFA, TLS và quản lý bí
mật giúp giảm đáng kể bề mặt tấn công so với việc mỗi ứng dụng tự lưu trữ thông
tin xác thực.

\section{Đánh giá mức độ tự động hóa}
\label{sec:autoscore}

Bảng~\ref{tab:autoscore} tổng hợp mức độ tự động hóa đạt được đối với từng hạng
mục của hệ thống.

\begin{table}[htbp]
\centering
\caption{Mức độ tự động hóa đạt được}
\label{tab:autoscore}
\small
\begin{tabular}{p{6.0cm} p{2.8cm} p{3.8cm}}
\toprule
\textbf{Hạng mục} & \textbf{Mức độ} & \textbf{Hình thức}\\
\midrule
Xây dựng image & Đạt & Pipeline CI (GitHub Actions)\\
Quản lý cấu hình định danh & Đạt & Realm JSON và keycloak-config-cli\\
Triển khai lên cụm & Đạt & Helm và manifest Kubernetes\\
Đồng bộ theo GitOps & Đạt & ArgoCD\\
Đo lường và giám sát & Đạt & Điểm cuối /metrics của Keycloak\\
Quản lý bí mật & Chưa & Biến môi trường\\
\bottomrule
\end{tabular}
\end{table}

\section{Đánh giá hiệu năng}
\label{sec:hieunang}

Để đánh giá sơ bộ khả năng đáp ứng của dịch vụ định danh, kịch bản sinh tải bằng
k6 (trình bày ở Phụ lục~\ref{lst:apx-k6}) được sử dụng để gửi yêu cầu cấp token
tới Keycloak với mức tải tăng dần. Kết quả ghi nhận cho thấy dịch vụ phản hồi ổn
định và không phát sinh lỗi trong suốt thời gian kiểm thử. Việc đo đạc định
lượng chi tiết được xác định là nội dung tiếp tục của giai đoạn phát triển.

\section{So sánh với phương án triển khai thủ công}
\label{sec:sosanh}

Bảng~\ref{tab:ss-trienkhai} so sánh phương án triển khai tự động đã xây dựng
với cách triển khai thủ công truyền thống.

\begin{table}[htbp]
\centering
\caption{So sánh triển khai thủ công và tự động}
\label{tab:ss-trienkhai}
\small
\begin{tabular}{p{4.2cm} p{4.6cm} p{4.6cm}}
\toprule
\textbf{Tiêu chí} & \textbf{Thủ công} & \textbf{Tự động (đã xây dựng)}\\
\midrule
Tính tái lập & Thấp, phụ thuộc người thao tác & Cao, tái lập từ Git\\
Thời gian triển khai & Lâu, dễ sai sót & Nhanh, nhất quán\\
Truy vết thay đổi & Khó & Có, qua lịch sử Git\\
Khả năng khôi phục & Chậm & Đồng bộ lại từ Git\\
Yêu cầu kỹ năng & Thấp & Cao hơn (đầu tư ban đầu)\\
\bottomrule
\end{tabular}
\end{table}

\section{Thảo luận}
\label{sec:thaoluan}

Bảng~\ref{tab:uunhuoc} tổng hợp ưu điểm và hạn chế của hệ thống đã xây dựng.

\begin{table}[htbp]
\centering
\caption{Ưu điểm và hạn chế của hệ thống}
\label{tab:uunhuoc}
\small
\begin{tabular}{p{6.2cm} p{6.6cm}}
\toprule
\textbf{Ưu điểm} & \textbf{Hạn chế}\\
\midrule
Xác thực tập trung, một nguồn định danh duy nhất & Phụ thuộc vào IdP (điểm tin cậy tập trung)\\
Cấu hình tái lập và truy vết qua Git & Đòi hỏi kỹ năng và công cụ hỗ trợ\\
Sẵn sàng chạy nhiều bản sao và tự phục hồi & Chưa được kiểm chứng ở quy mô lớn\\
Bảo mật theo lớp (TLS, MFA, RBAC) & Quản lý bí mật còn ở mức cơ bản\\
\bottomrule
\end{tabular}
\end{table}

Việc kết hợp xác thực tập trung với triển khai tự động cho thấy đây là hướng
tiếp cận phù hợp: dịch vụ định danh vừa được chuẩn hóa về giao thức, vừa được
vận hành có kiểm soát và tái lập. Các hạn chế nêu trên là cơ sở cho những nội
dung phát triển tiếp theo được trình bày ở phần Kết luận.

\section{Kết chương}

Chương này đã trình bày kết quả kiểm thử chức năng, kết quả đánh giá bảo mật
theo mô hình STRIDE, mức độ tự động hóa đạt được, đánh giá hiệu năng sơ bộ và so
sánh với phương án triển khai thủ công.
